\chapter{The KV cache and its band-aids}
\label{ch:kv-cache}

\newthought{FlashAttention} fixes the activation-memory bill of the
\(n \times n\) attention matrix during a single forward pass. It does not
fix the per-token KV cache that grows during autoregressive generation.
This chapter is the catalogue of partial fixes for that second bill.

\section{Why caching is necessary}

Recall the autoregressive loop. To generate token \(t+1\), the model needs
to attend over tokens \(1, \ldots, t\). To compute attention, it needs
their keys and values. Recomputing \(K_{1:t}\) and \(V_{1:t}\) from scratch
at every step would make generation \(\bigO(n^2 d)\) per token, i.e.\
\(\bigO(n^3)\) for the whole sequence. Caching makes it \(\bigO(n d)\) per
token, i.e.\ \(\bigO(n^2)\) overall — back to the same asymptote, but with
a much friendlier constant.

So at every layer, every head, the model maintains running tensors:
\begin{equation}
  K^{(\ell, h)} \in \R^{n \times \dhead}, \quad
  V^{(\ell, h)} \in \R^{n \times \dhead}.
\end{equation}
Per token of generation, the bytes added are
\begin{equation}
  2 \cdot L \cdot \heads \cdot \dhead \cdot s
\end{equation}
where \(s\) is bytes per scalar. Earlier we worked the numbers: roughly
0.5 MB per token for a 7B-class model, half a terabyte for 1M tokens.

Everything that follows is an attempt to reduce that constant.

\section{Multi-Query Attention (MQA)}

\citet{shazeer2019mqa} proposed the simplest fix. Use \(\heads\) different
\emph{queries}, but share a \emph{single} key and value tensor across all
heads:
\begin{equation}
  Q_h \in \R^{n \times \dhead}, \quad K \in \R^{n \times \dhead}, \quad V \in \R^{n \times \dhead}.
\end{equation}
The KV cache shrinks by a factor of \(\heads\). For 32 heads, the per-token
cache drops from 0.5 MB to ${\sim}16$ KB.\sidenote{Quality typically
degrades a few percentage points on standard benchmarks; the interpretation
is that the 32 heads were doing genuinely different work and collapsing
their keys throws some of that diversity away.}

\section{Grouped-Query Attention (GQA)}

\citet{ainslie2023gqa} compromise: split the heads into \(g\) groups, share
keys and values within each group. \(g = 1\) is MQA; \(g = \heads\) is
standard MHA. Llama 2 70B and Llama 3 use GQA with \(g = 8\). The cache
shrinks by \(\heads / g\); the quality hit is much smaller than MQA's.

\begin{center}
\begin{tabular}{lll}
\toprule
Scheme & KV cache scaling & Quality cost \\
\midrule
MHA  (\(g = \heads\)) & full                 & none (baseline) \\
GQA  (\(g = 8\))      & \(\heads/8\) smaller & negligible in practice \\
MQA  (\(g = 1\))      & \(\heads\)$\times$ smaller & 1--3 pp on most benchmarks \\
\bottomrule
\end{tabular}
\end{center}

\section{Paged KV (vLLM)}

\citet{kwon2023vllm} treat the KV cache the way an OS treats memory: split
each sequence's cache into fixed-size pages and maintain a page table
mapping logical token positions to physical blocks. Multiple sequences
share the underlying memory pool; pages can be added on demand and freed
when the sequence ends or is preempted.

The win is throughput, not memory per se: paged KV lets you serve many
concurrent requests on the same GPU without fragmenting memory or
over-provisioning. It is the standard inference layer for Llama-class
models in 2026 (vLLM, TGI, TensorRT-LLM).

\section{Sliding window attention (SWA)}

Mistral's signature trick \citep{jiang2023mistral7b}. Instead of attending
to all previous tokens, each token attends only to the last \(w\) (typically
\(w = 4096\)). The KV cache size becomes \(\bigO(w)\) instead of \(\bigO(n)\).
Mistral 7B layered SWA on top of the standard transformer and combined it
with the trick that successive layers' windows compose, so the
\emph{effective} receptive field at the top layer is
\(L \cdot w\).\sidenote{For \(L = 32\), \(w = 4096\), effective receptive
field is roughly 130K tokens. Less than a true 130K-token attention,
because the receptive field is convolutional in shape — distant tokens'
information has to be relayed through intermediate layers — but a
real-world win.}

The SWA trade-off is precisely the one we will revisit in Part~V: cheap,
but the routing is fixed by position, so the model literally cannot see
information outside the window without help.

\section{RoPE, YaRN, and position extrapolation}

Rotary position embeddings (RoPE, \citealp{su2021rope}) encode position by
rotating the query and key vectors in 2D subspaces by an angle proportional
to position. RoPE has the elegant property that the dot product
\(q_i \cdot k_j\) only depends on \(j - i\), so it composes naturally with
relative attention.

If you train at sequence length \(n_{\text{train}}\) and want to deploy at
\(n_{\text{infer}} \gg n_{\text{train}}\), naive RoPE breaks because the
rotations land in unfamiliar regions of the angle space. YaRN
\citep{peng2023yarn} and a family of related tricks rescale the rotation
frequencies so that the model can extrapolate to longer contexts than it
was trained on, with much less degradation.

These are not algorithmic changes to attention itself — they are clever
position-encoding tweaks that let you stretch the trained model's effective
range. They do not change the asymptotic cost, and they do not, by
themselves, recover the lost reasoning quality at extreme \(n\).

\section{The honest summary}

Each of MQA, GQA, paged KV, SWA, and RoPE/YaRN is a real win. Together
they are the reason production transformers can serve 128K--1M token
contexts at all in 2026. None of them changes the \(\bigO(n^2)\) compute
or the fundamental fact that the KV cache grows linearly per layer.

\keyidea{Every mitigation in this chapter is a band-aid on an
\(\bigO(n^2)\) wound. The patient survives. The disease is still there.
At sufficiently large \(n\), or sufficiently demanding quality bars, you
have to change the algorithm — which is what the rest of the book is
about.}

\begin{table}[htbp]
\centering
\small
\begin{tabular}{@{}lrrr@{}}
\toprule
KV layout & KV total (\(\mathrm{GiB}\), \textsc{fp16}) & H100 \(80\,\mathrm{GiB}\) req.\ / GPU & B200 \(192\,\mathrm{GiB}\) req.\ / GPU \\
\midrule
MHA (64 KV heads) & \(\approx 2{,}441\) & 0 & 0 \\
GQA-8 (Llama-3 70B) & \(\approx 305\) & 0 & 0 \\
MQA (1 KV head) & \(\approx 38.1\) & \(\le 2\)\textsuperscript{*} & \(\le 5\)\textsuperscript{*} \\
\bottomrule
\end{tabular}
\caption{KV cache only for one concurrent request at \(n=10^{6}\) tokens,
\(L{=}80\), \(\dhead{=}128\), \(\heads{=}64\). Numbers scale as
\(2 L n\, n_{\text{kv}}\, \dhead\) half-precision bytes.
\textsuperscript{*}Concurrent count is \(\lfloor \mathrm{GPU\ memory} / \mathrm{KV}\rfloor\) ignoring weights, activations, and fragmentation---real serving fits fewer.}
\label{tab:kv-70b-1m}
\end{table}

\noindent
GQA-8 is why production 70B models are servable at 128k--1M on \emph{pods} of
GPUs rather than on a single device: KV alone already exceeds one H100 for
million-token \emph{single-sequence} residency. MQA trades representational
freedom for capacity; Jamba-class hybrids (Chapter~\ref{ch:hybrids}) attack the
same wall by mixing in linear-time blocks.

\paragraph{PagedAttention = virtual memory for KV blocks.}
\citet{kwon2023vllm} treat KV tensors as fixed-size \emph{pages} (blocks of
tokens) allocated from a GPU-wide pool. A sequence's logical token order maps
through a small indirection table to physical pages, just as virtual addresses
map to physical frames in an OS.
Non-contiguous physical placement avoids the classic ``reserve one giant
rectangle of memory for the worst-case prompt length'' fragmentation problem:
different sequences grow by allocating new pages; short sequences do not waste
holes sized for someone else's 1M-token tail.
Parallel sampling from a shared prompt is handled copy-on-write: child beams
share read-only prompt pages until a branch diverges, then duplicate only the
suffix pages that differ---analogous to \texttt{fork} + COW in Unix kernels.
The net effect is higher batching on the same hardware without changing model
math.

\paragraph{RoPE, NTK-aware scaling, YaRN, and position extrapolation.}
RoPE \citep{su2021rope} injects position by rotating \(Q\) and \(K\) in paired
planes with frequencies \(\theta_k\) spaced across a geometric ladder.
Short training lengths mean high-frequency components complete many cycles over
the training horizon; when inference extends far beyond that horizon, those
components enter \emph{unseen} phases, so attention behaves as if it left the
training manifold.
\emph{NTK-aware} rescaling (popularised in community writeups by kaiokendev,
bloc97, and others) borrows the Neural Tangent Kernel intuition: shrink effective
frequencies so rotations advance more slowly at long positions, trading some
fine local discrimination for stability.
YaRN \citep{peng2023yarn} generalises this idea with learned or hand-tuned
stretching of the frequency ladder (sometimes mixing ``attention sinks'' and
temperature-like scaling) so checkpoints trained at \(L\) tokens reach \(2L\),
\(8L\), or more with less perplexity blow-up.
\emph{Position interpolation} (linearly rescaling position indices) is the
simplest baseline in the same family.
None of these tricks changes \(\Theta(n^{2})\) attention cost or KV linear
growth; they only stretch the \emph{usable} portion of the context window after
training. The honest recipe for a new length regime remains: train (or at least
adapt) there, then evaluate on long benchmarks such as RULER/MRCR, not just
PPL on held-out slices.
